% ============================================================
%  CHAPTER 8 — Navigation
% ============================================================
\chapter{Navigation}
\label{ch:navigation}

\section{Overview}

Navigation provides the robot with the ability to move between named locations in
its environment — for example, from a home position to the table where objects are
placed, or to a person waiting to receive an object.
The navigation subsystem is built on the ROS \texttt{move\_base} stack, which
integrates global and local path planning, obstacle avoidance, and recovery behaviours.

In the current deployment, navigation is implemented and tested but not yet integrated
into the main VLA agent loop.
The agent operates in a fixed position, relying on the arm's reach and head scanning
to interact with objects.
Navigation integration is identified as a key future improvement (Chapter~\ref{ch:conclusion}).

\section{Navigation Stack}

\subsection{ROS move\_base}

The \texttt{move\_base} node implements a complete mobile base navigation solution:

\begin{itemize}[noitemsep]
  \item \textbf{Global planner}: computes an initial path from current pose to goal
        using Dijkstra's algorithm on a static map.
  \item \textbf{Local planner}: the Dynamic Window Approach (DWA) continuously adapts
        the path in response to sensor data, enabling real-time obstacle avoidance.
  \item \textbf{Costmaps}: inflation layers around obstacles discourage the robot from
        passing too close to walls and furniture.
  \item \textbf{Recovery behaviours}: if the robot becomes stuck, it executes a
        sequence of escalating recovery actions — rotating in place, clearing the
        costmap, and backing up.
\end{itemize}

Navigation goals are sent via the \path{/move_base} action server
(\texttt{move\_base\_msgs/Move\-BaseAction}).

\subsection{Navigation Client}

The \texttt{NavigationClient} class (\texttt{navigation\_client.py}) wraps the
\texttt{move\_base} action client with a named-location dictionary:

\begin{lstlisting}[language=Python, caption={Navigation to a named location.}]
NAMED_LOCATIONS = {
    'home':  {'x': 0.0,  'y': 0.0,  'theta': 0.0},
    'table': {'x': 1.5,  'y': 0.3,  'theta': 0.0},
    'person':{'x': 0.8,  'y':-0.5,  'theta': 1.57},
}

def go_to(self, location_name, timeout=60):
    loc = NAMED_LOCATIONS[location_name]
    goal = MoveBaseGoal()
    goal.target_pose.header.frame_id = "map"
    goal.target_pose.pose.position.x = loc['x']
    goal.target_pose.pose.position.y = loc['y']
    # Convert theta to quaternion
    q = quaternion_from_euler(0, 0, loc['theta'])
    goal.target_pose.pose.orientation.z = q[2]
    goal.target_pose.pose.orientation.w = q[3]
    self.client.send_goal(goal)
    return self.client.wait_for_result(rospy.Duration(timeout))
\end{lstlisting}

\section{Location Management}

\subsection{Saving Locations}

The \texttt{save\_location.py} script allows the operator to drive the robot to a
desired position and save it with a name:

\begin{lstlisting}[language=Python, caption={Saving current robot pose as named location.}]
# Get current pose from odometry
odom = rospy.wait_for_message('/mobile_base_controller/odom', Odometry)
x = odom.pose.pose.position.x
y = odom.pose.pose.position.y

# Save to YAML
locations[name] = {'x': x, 'y': y, 'theta': theta}
with open('locations.yaml', 'w') as f:
    yaml.dump(locations, f)
\end{lstlisting}

Saved locations persist across robot restarts and are loaded at agent startup.

\subsection{Robot Pose Monitoring}

\texttt{check\_robot\_pose.py} continuously prints the robot's current position in
the map frame, useful during manual positioning and debugging:

\begin{lstlisting}[language=bash]
# Output example:
[Pose] x=1.523 y=0.312 theta=0.015 rad (0.9 deg)
\end{lstlisting}

\section{Autonomous Mapping}

\texttt{autonomous\_mapper.py} uses the GMapping SLAM algorithm to build a 2-D
occupancy grid map of the environment while the robot navigates autonomously.
The robot is teleoperated or follows a pre-defined exploration trajectory while
the laser scanner data is used to incrementally build the map.
The resulting map is saved in YAML/PGM format and loaded by \texttt{move\_base} for
subsequent navigation sessions.

\section{Recovery Behaviours}

The \texttt{unstuck\_robot.py} script implements manual recovery when the robot's
navigation stack declares it stuck (all recovery behaviours exhausted):

\begin{enumerate}[noitemsep]
  \item Publish a zero-velocity command to stop any ongoing motion.
  \item Clear the local costmap via the \texttt{/move\_base/clear\_costmaps} service.
  \item Rotate the base $\pm$0.5 rad in place.
  \item Retry the original navigation goal.
\end{enumerate}

\section{Base Velocities for Simple Motions}

For cases where full navigation planning is not needed — such as rotating to face
a detected object or backing away from an obstacle — \texttt{quick\_move\_base.py}
and \texttt{turn\_and\_go.py} publish directly to \texttt{/cmd\_vel}:

\begin{lstlisting}[language=Python, caption={Publishing a rotation command.}]
cmd = Twist()
cmd.angular.z = 0.3   # rad/s counterclockwise
pub.publish(cmd)
rospy.sleep(duration)
cmd.angular.z = 0.0
pub.publish(cmd)      # stop
\end{lstlisting}

\section{Integration with the Agent}

Navigation is scaffolded in the agent via the \texttt{NavigateToSkill} in
\texttt{skills/navigate\_to.py}.
The skill checks affordance by verifying that \texttt{move\_base} is available
(1-second timeout on the action server) and that the requested location name exists
in the location dictionary.
On execution, it calls \texttt{NavigationClient.go\_to()} with a 60-second timeout.

The skill is not currently loaded into the agent's skill library by default because:
\begin{itemize}[noitemsep]
  \item Navigation requires a pre-built map specific to the deployment environment.
  \item Object positions (table, handover point) must be consistent with the map frame.
  \item In the test environment, the robot can reach all relevant objects from
        a single fixed position.
\end{itemize}

To enable navigation, load the skill in \texttt{EmbodiedAgent.\_load\_skills()}:

\begin{lstlisting}[language=Python, caption={Enabling navigation in the agent.}]
from skills.navigate_to import NavigateToSkill
self.skills['navigate_to'] = NavigateToSkill()
\end{lstlisting}

and add a corresponding description in the VLM skill list.
